\documentclass[a4paper]{article}
\usepackage[utf8]{inputenc}
\usepackage{listings}
\usepackage{parcolumns}
\usepackage{hyperref}
\usepackage{graphicx}
\usepackage{float}
\usepackage{pifont}
\newcommand{\cmark}{\ding{51}}
\newcommand{\xmark}{\ding{55}}
\newcommand{\umark}{\xmark/\cmark}
\newcommand{\dmark}{\cmark/\xmark}
\usepackage{amsmath}
\usepackage{tikz}
\usepackage{makecell}
\usetikzlibrary{shapes.misc, arrows.meta}
\usepackage{array}
\newcolumntype{C}{>$c<$}
\usepackage{color}
\usepackage{fancyhdr}
\usepackage{geometry}
\usepackage[T1]{fontenc}
\usepackage[utf8]{inputenc}
\usepackage{lastpage}
\usepackage{listings}
\usepackage{enumitem}
\usepackage{tabu}
\usepackage[flushleft]{threeparttable}
\usepackage{ulem}
\usepackage{amsmath}
\usepackage{pdfpages}
\usepackage{multicol}
\usepackage{tabu}
\usepackage{parskip}
\usepackage{qtree}
\usepackage{tcolorbox}
\usepackage{scalerel,amssymb}
\def\mcirc{\mathbin{\scalerel*{\circ}{j}}}
\def\msquare{\mathord{\scalerel*{\Box}{gX}}}
\newcommand{\code}[1]{\begin{tcolorbox} \texttt{#1} \end{tcolorbox}}
\def\doubleunderline#1{\underline{\underline{#1}}}
\begin{document}

% COURSE NAME
\title{02561 - Computer Graphics}
\author{
\includegraphics[width=0.15\textwidth]{images/dtu.eps}~\\[1cm]
    DTU - Technical University of Denmark
    \\[0.5cm]
    Date of submission: \today
    \\
}
\date{} % blind date
\color{black}
\maketitle
\begin{center}
{ \huge \bfseries Virtualizing physical object movements}\\

\vspace{.25cm}
Daniel F. Hauge \texttt{(s201186)}\\

% \vspace{.25cm}
% Name \texttt{(s12345678)}\\

\vspace{.25cm}
\end{center}

\begin{abstract}
This project aims to virtualize and display a physical object with it's physical movements on screen. 
The project consists of an implementation which succesfully captures 3 degrees of freedom and uses the movements to display on screen.
\\
{ \begin{center} \bfseries The implementation is hosted on: \href{https://grafik.feveile-hauge.dk}{grafik.feveile-hauge.dk}. \end{center}}
\end{abstract}

\medskip
\newpage

\section{Introduction}
With ever growing computing power at disposal, more opportunities for cool technology emerges. One such technology that has just recently picked up speed, 
is virtual reality and augmented reality. I postulate that virtual reality is by no means a matured complete technology, as only two human senses has convincingly been implemented.
The cutting edge VR technology still mostly require all movements to happen in actual reality, ie. movements happen in reality and are captured, virtualized and rendered in VR.
This project will try to tackle the problem of capturing movements to display physical movements in a virtual environment.

\subsection{Problem statement}
How can physical movement be captured and displayed in a virtual environment like on a HTML canvas.

\subsection{Motivation \& Usages}
As previusly mentioned, VR technology still require bodily movements to happen in reality.
So capturing physical movements is currently still relevant within VR technology.
Just like VR, augmented reality also benefit from being able to capture physical movements to be used for displaying graphics.

Capturing physical movements is also relevant in film making.
Movies like Lord of the Rings and Avatar, have used motion capture technology to great success. 

Apart from being used for rendering graphics,
capture of movements can be used in new computer interfaces, 
like some sort of gesture controller.  
\section{Method}
There are challenges the problem at hand etc.

\subsection{Capture physical movements}
Sensors. Magnetic \& Accelerometors Light? etc.
Choosing Magnetic and accelerometors, as it can be self contained etc.

The phone has these sensors and can be used as a sort of controller for capturing movements.

\subsection{Data transfer}
Wireless transfer to allow for 100\% freedom in movements. Bluetooth, wifi.
Requirement: 100\% freedom in movements. The object should only be under constraints as in the real world.
A phone is still a good candidat for this requirement, as it is wireless. 

Wifi is used picked, as it is adequately fast and reliable.

\subsection{Display}
Using a triangle mesh to instantiate vertex. 
We can use affine transformation with a rotation and translation matrix to move the object in virtual space.
\section{Implementation}

\subsection{Capturing physical movements}
\subsubsection{Orientation}
To make an implementation that would work for most devices, a widely available API should be used. There is support for \textbf{AbsoluteOrientationSensor} API in 
most modern browsers. This sensor provides an API for capturing orientation of a mobile device. The API outputs quaternion, which is very usefull for the currently choosen method.
To use the phone with the web api for capturing movements, a new page on the website is created for the phone to visit.
When visiting the page from a phone, the sensor data is captured and sent to another device which is visiting the display page.
The sensor is initialized with frequency and reference frame, and a event listener can be added before starting.

The orientation sensor is setup and started with the following code snippet:
\code{orientationSensor = new AbsoluteOrientationSensor(\{ frequency: 60, referenceFrame: 'device' \}); \\
orientationSensor.addEventListener('reading', ev => \{ \\
\hspace*{10mm} const x = orientationSensor.quaternion[0]; \\
\hspace*{10mm} const y = orientationSensor.quaternion[1]; \\
\hspace*{10mm} const z = orientationSensor.quaternion[2]; \\
\hspace*{10mm} const w = orientationSensor.quaternion[3]; \\
\hspace*{10mm} console.log(x, y, z, w); \\
\});\\
orientationSensor.start();
}

On each read, the sensor object has an updated property called \textbf{quaternion}, which is the last captured physical rotational movements.

\subsubsection{Displacement}
Just like orientation, displacement is captured with a web api which is supported in most modern browsers. 
The page that captures phone data is extended to capture acceleration aswell. 
The \textbf{LinearAccelerationSensor} is used for capturing acceleration.
Just like the orientation sensor, it is initialized with a frequency and an event listener can be added.
Each readings contain an acceleration along x, y and z axis.

\code{accelerationSensor = new LinearAccelerationSensor(\{frequency: 60\});\\
accelerationSensor.addEventListener('reading', () => \{\\
\hspace*{10mm} const x = accelerationSensor.x;\\
\hspace*{10mm} const y = accelerationSensor.y;\\
\hspace*{10mm} const z = accelerationSensor.z;\\
\hspace*{10mm} console.log(x, y, z, w);\\
\});\\
accelerationSensor.start();
}

On each read, the sensor object has updated x, y and z properties, which correspond to the accelerations in x, y and z axis.

\subsubsection{Support}
An important implementation detail for using sensor web api's, is that it is required to access the website through a secure connection.
In other words, HTTPS is required and regular HTTP is not adequate.
Another thing is, that if the device does not have adequate sensors, the API's will not be defined in the javascript environment.
Therefor error handling has also been implemented to let a user know if sensor is unavailable and more. 

\subsection{Data transfer}
The data transfer has to happen between 2 clients of the website.
One client is the phone that is visiting the page which setup capture sensors.
Another client is a device that can display the virtual environment.
The implementation is done with websockets. I found a golang and a javascript library from \href{https://socket.io/}{socket.io} that implements websockets.

Using the library, dedicated connection logic is implemented, such that a phone can connect to another client.
There is quite a bit more to how the data transfer implementation is done, but to keep it simple, only the actual transfering of data is covered.
See \textbf{websocket.js} and \textbf{main.go} for full websocket code. 

When a connection is established, the phone can then emit it's readings to the client via websockets.
The following statement is emiting events with orientation sensor readings to the server.

\code{ProjectSockets.Socket.emit("event", \{ id: connectionId, type: "orientation", value: JSON.stringify(orientationSensor.quaternion)\});}

The website server is extended with a websocket server handler:
Listeners for specific channels can then be added. 
The following listener is adding functionality for the server to receive the phone's events and broadcast them to the connected display client.
\code{server.On("event", func(c *gosocketio.Channel, msg Event) string \{ \\
   \hspace*{10mm} if msg.EventType == "connection" \{ \\ 
   \hspace*{20mm}     if observer[msg.SessionId] || !listener[msg.SessionId] \{ \\
   \hspace*{30mm}     return "BAD" \\
   \hspace*{20mm}   \} else \{ \\ 
   \hspace*{30mm}      observer[msg.SessionId] = true \\
   \hspace*{30mm}   connections[c.Id()] = msg.SessionId \\ 
   \hspace*{20mm}    \} \\
   \hspace*{10mm}   \} \\
   \hspace*{10mm} c.BroadcastTo(msg.SessionId, msg.EventType, msg.Value) \\ 
   \hspace*{10mm}    return "OK" \\
\})
}

The following code will setup a listener on the display client that listens for events from the phone through the server. The event contain the readings from the phone.

\code{ProjectSockets.Socket.on("orientation", orientationArray => \{ \\
    \hspace*{10mm}  let quaternion = JSON.parse(orientationArray); \\
    \hspace*{10mm}   if (ProjectSockets.OnOrientationRead != null) \{ \\
        \hspace*{20mm}      ProjectSockets.OnOrientationRead(quaternion); \\
        \hspace*{10mm}   \} \\
        \}); 
}

This listener calls a function \textbf{OnOrientationRead} on every received readings from the phone.
Similar listeners are made with acceleration, connection, disconnection and alignment.
The function can be assigned to let the display device react to new orientations et cetera.


\subsection{Display}
To display the environment, first a virtual environment with a static object is implemented.
Using the same methods and procedures as in worksheet 5 and 10, an object is loaded with a .OBJ file, and a quad is loaded with a texture for looks.

To rotate the object based on the phones readings, a variable holding a quaternion is declared and updated by each readings from the phone.
The following code assigns the \textbf{OnOrientationRead} function to set a reference orientation and update the quaternion variable with respect to the reference.
\code{ProjectSockets.OnOrientationRead = ori => \{ \\
\hspace*{10mm} const newOri = [ori[0], ori[1], ori[2], ori[3]] \\
\hspace*{10mm} if (Project.refSet == undefined) \{ \\
\hspace*{20mm} Project.refSet = true; \\
\hspace*{20mm} Project.q\_rot\_ref.set(newOri); \\
\hspace*{10mm} \} \\
\hspace*{10mm} Project.q\_rot.set(newOri); \\
\hspace*{10mm} Project.q\_rot.multiply(Project.q\_rot\_ref); \\
\}
}

The quaternion.js contain a function for calculating the matrix, this is used when drawing the phone:

\code{const rot = new Matrix4(); \\
const mat4 = Project.q\_rot.get\_mat4(); \\
const flattened = flatten(mat4); \\
rot.set({elements: flattened}); \\
Project.g\_modelMatrix = Project.g\_modelMatrix.multiply(rot);
}

The phone readings are now able to rotate a model matrix which is then used for drawing an object in a virtual environment.
There are more implementation details which involves perspective alignment with the screen and more, these will not be covered to keep it succinct.
See \textbf{project.js} for full code on display/drawing.

To draw the displacement, an attempt with similar approach to orientation is done.
The following code is assigning the \textbf{OnAccelerationRead} to calculate a new total displacement vector.

\code{const Displacement = \{x:0, y:0, z:0\}; \\
const Velocity = \{x:0, y:0, z:0\}; \\
const DeltaT = 1/15; \\
ProjectSockets.OnAccelerationRead = Acceleration => \{ \\
\hspace*{10mm}   const vX = Velocity.x + Acceleration.x * DeltaT \\
\hspace*{10mm}  const vY = Velocity.y + Acceleration.y * DeltaT \\
\hspace*{10mm}  const vZ = Velocity.z + Acceleration.z * DeltaT \\
\hspace*{10mm}   Displacement.x = 0.5 * (vX + Velocity.x) * DeltaT \\
\hspace*{10mm}  Displacement.y = 0.5 * (vY + Velocity.y) * DeltaT \\
\hspace*{10mm}   Displacement.z = 0.5 * (vZ + Velocity.z) * DeltaT \\
\hspace*{10mm}   Velocity.x = vX; \\
\hspace*{10mm}   Velocity.y = vY; \\
\hspace*{10mm}   Velocity.z = vZ; \\
\hspace*{10mm}   Project.PD.x =  Displacement.x; \\
\hspace*{10mm}   Project.PD.y =  Displacement.y; \\
\hspace*{10mm}   Project.PD.z =  Displacement.z;\\
    \} \\
}

A translation matrix is then multiplied to the model matrix:

\code{Project.g\_modelMatrix.translate(\\ \hspace*{10mm} Project.Phone\_X +  Project.PD.x, \\ \hspace*{10mm} Project.Phone\_Y + Project.PD.y, \\ \hspace*{10mm} Project.Phone\_Z + Project.PD.z);}

Unfortunately, the implementation for displaying displacement is not working as intended.
It seems as through acceleration is not accurate, as moving the phone from one stationary position to another, the velocity is not zero when the phone is done moving.
The acceleration or deacceleration reads must have been larger to produce a non zero velocity.
A different method could perhaps be better at implementing the solution.
\section{Solution}
- The solution is this and that.
With video.

\subsection{Trying the solution}
- Instructions.
\section{Conclusion}
- Sensors from a phone can be used to capture and virtualize movements.
- The resulting solution is not super usefull but demonstrates potential for other applications.
- Unfortunately the solutions does not support displacement, which would have been cool.

% Sensors from modern smartphones can be used to capture and virtualize movements.


% Credits: https://open3dmodel.com/3d-models/samsung-galaxy-s20_477030.html

\end{document}







